\section{The least orbit dimension and uniform small balls}\label{sec:minorbits58}
The geometric hypothesis needed for a stochastic converse concerns every
conditional-centroid direction, not only the prescribed target orbit.
For a fixed finite-dimensional orthogonal representation of a compact
connected Lie group, this hypothesis follows from a rank invariant.
Throughout this section Haar measure is normalized, and the Lie algebra
has a fixed invariant inner product. No constants uniform in the dimension
or in the representation are asserted.

\begin{theorem}[Uniform caps from the least orbit dimension]
\label{thm:minorbit58}
Let $G$ be a compact connected Lie group and $R:G\to O(E)$ a continuous
finite-dimensional representation with $E^G=\{0\}$ and $E\ne\{0\}$.
Write $D=\dim G$ and choose a basis $X_1,\ldots,X_D$ of its Lie algebra,
using the same symbols for their infinitesimal actions on $E$. Set
\begin{equation}\label{eq:minrank58}
 p_* =\min_{\norm u=1}\operatorname{rank}
       [X_1u\ \cdots\ X_Du]
     =\min_{\norm u=1}\dim R(G)u.
\end{equation}
Then $p_*\ge1$, and there is $C<\infty$ such that
\begin{equation}\label{eq:mincap58}
 \sup_{\norm u=1,\ v\in E}
 \mu\{g:\norm{R(g)u-v}<h\}\le C h^{p_*}
                 \qquad(0<h\le1).
\end{equation}
The exponent is optimal as a uniform exponent: it cannot be replaced by
any larger exponent with a constant independent of $u,v,h$.
\end{theorem}
\begin{proof}
The tangent space to the orbit at $u$ is the span of the $X_i u$.
An orbit of dimension zero is a singleton, since $G$ is connected; its
point would be fixed. This proves $p_*\ge1$. Put $r=p_*$ and fix
orthonormal coordinates on $E$.

For each unit $u$, some $r$-row, $r$-column minor of
$[X_1u\ \cdots\ X_Du]$ is nonzero. The maximum of the absolute values
of these finitely many minors is a continuous positive function on the
unit sphere. Its minimum is therefore positive. The entries are uniformly
bounded. Consequently there is $\sigma>0$ such that for every unit $u$
one can choose $r$ input indices and $r$ output coordinates whose derivative
matrix $A_u$ has least singular value at least $\sigma$.

For each of the finitely many permutations of the chosen Lie algebra
basis consider the product chart
\[
 \phi(x,y)=\exp(x_1X_{i_1})\cdots\exp(x_rX_{i_r})
            \exp(y_1X_{i_{r+1}})\cdots\exp(y_{D-r}X_{i_D}).
\]
All these charts are diffeomorphisms on one sufficiently small cube about
zero. Their images contain a common identity neighborhood $W$, and their
Haar densities are bounded above by a common constant. By shrinking the
cube once more, uniform continuity of the derivatives on the unit sphere
and the finite list of charts ensures the following. For every unit $u$,
choose the minor above and let $\pi_u:E\to\R^r$ be its coordinate
projection. For all $(x,y)$ in the cube,
\begin{equation}\label{eq:submersion58}
 \left\|\partial_x\bigl(\pi_u R(\phi(x,y))u\bigr)-A_u\right\|
                                                    \le\sigma/2.
\end{equation}
For fixed $y$, the map $f_{u,y}(x)=\pi_u R(\phi(x,y))u$ is injective.
Indeed integration of the derivative on the segment between $x$ and
$x'$ gives
\[
 \norm{f_{u,y}(x)-f_{u,y}(x')}
       \ge\norm{A_u(x-x')}-\tfrac12\sigma\norm{x-x'}
       \ge\tfrac12\sigma\norm{x-x'}.
\]
Every singular value of its derivative is at least $\sigma/2$.
The change-of-variables formula therefore bounds the $r$-dimensional
volume of the inverse image of a ball of radius $h$ by
$\omega_r(2/\sigma)^r h^r$. A ball in $E$ projects into a ball of the
same radius in these coordinates. Integrating in $y$ and using the bounded
Haar density proves the required estimate on this chart, uniformly in
$u$ and in the center $v$.

Choose finitely many right translates $Wg_j$ covering $G$. On such a
translate apply the preceding argument to the unit vector $R(g_j)u$.
For its selected minor, use the corresponding permuted chart; that chart
contains $W$. Haar invariance and a finite sum give \eqref{eq:mincap58}
for sufficiently small $h$. Enlarging $C$ covers the remaining
$h\le1$. This proof does not require a spectral decomposition of the
arbitrary center $v$.

Finally choose $u_*$ whose orbit has dimension $r$. The orbit is the
smooth embedded compact homogeneous space $G/G_{u_*}$. The pushforward
of Haar measure is its invariant probability measure, with positive
smooth density in local coordinates. At $u_*$, sufficiently small
ambient chordal balls have measure bounded below by $c h^r$ for some
$c>0$. Thus no larger uniform exponent is possible.
\end{proof}

\begin{proposition}[An algebraic rank certificate]\label{prop:rankcompute58}
Suppose the infinitesimal matrices $X_i$ are supplied with exact real
algebraic entries, and are known to specify the representation in
Theorem~\ref{thm:minorbit58}. Then $p_*$ is computable by real quantifier
elimination. A positive algebraic lower bound for the largest
$p_*\times p_*$ minor on the unit sphere is computable as well. If a unit
seed $z_0$ is algebraic, its orbit dimension $q$ is the rank of
$[X_1z_0\ \cdots\ X_Dz_0]$.
\end{proposition}
\begin{proof}
For each integer $r$ ask whether a unit vector exists on which all
$(r+1)\times(r+1)$ minors vanish. These are finitely many polynomial
conditions over real algebraic numbers; the first affirmative answer is
$p_*$. With $r=p_*$, the sum of squares of all $r$-minors has a positive
minimum on the unit sphere, by the first paragraph of the preceding
proof. Quantifier elimination computes that minimum as a real algebraic
number with isolating data. If there are $M$ minors, their largest
absolute value is at least $\sqrt{\gamma/M}$, where $\gamma$ is that minimum. Algebraic rank at the supplied seed gives the last assertion.
\end{proof}
The proposition computes the rank invariant from \emph{supplied}
infinitesimal matrices. It neither certifies that arbitrary input matrices
integrate to a compact group nor computes a random-walk spectral gap.

\begin{corollary}[Uniform width exponents without a separate cap hypothesis]
\label{cor:allrepresentations58}
In the orbit realization model of Theorems~\ref{thm:uniformorbit57} and~\ref{thm:sharpoccupation61}, assume
the declared spectral gap, but replace its small-ball hypothesis by
$E^G=\{0\}$. Let $p_*$ be \eqref{eq:minrank58}, and let $q$ be the real
dimension of the unit target orbit spanning $E$. For fixed
$0<\rho\le1$ and $0\le\varepsilon<\rho$,
\[
 W_{N,\varepsilon}\ge cN^{p_*/2}.
\]
For $0<\rho<1$ the exact upper bound is $C(N+1)^{q/2}$; for $\rho=1$
it holds at each fixed positive error. Every fixed width has the
occupation bound of Theorems~\ref{thm:uniformorbit57} and~\ref{thm:sharpoccupation61}. Error is measured in
the Euclidean norm defining the legal orbit hull.
\end{corollary}
\begin{proof}
Apply Theorem~\ref{thm:minorbit58} to obtain its uniform cap estimate with
exponent $p_*$. The lower and upper arguments of
Theorems~\ref{thm:uniformorbit57} and~\ref{thm:sharpoccupation61} then apply without alteration.
\end{proof}
The group dimension $D$, least orbit dimension $p_*$, and target orbit
dimension $q$ need not coincide. For the standard sphere representation
$p_*=q=d-1$. For conjugation on traceless Hermitian $d\times d$ matrices,
an orbit with eigenvalue multiplicities $m_1,\ldots,m_b$ has dimension
$d^2-\sum_i m_i^2$. A nonzero traceless matrix has $b\ge2$, and the
largest possible sum is $(d-1)^2+1$. Thus $p_*=2(d-1)$, attained by a
centered rank-one projection. The explicit Grassmannian estimates below
also provide this exponent; they are retained as a direct geometric
proof with stated dimension-dependent constants.

The converse now uses Theorem~\ref{thm:sharpoccupation61}, rather than
paying the pointwise-mixing block length at each narrow cut. Thus no
logarithm separates the bounds when $p_*=q$. A supplied algebraic
infinitesimal representation makes $p_*$ computable; this statement
does not compute a spectral gap or assert dimension-uniform constants.
